\documentclass[11pt,a4paper]{article}
\usepackage[spanish]{babel}
\usepackage[utf8]{inputenc}
\usepackage[T1]{fontenc}
\usepackage{lmodern,geometry,xcolor,xurl,hyperref,enumitem,booktabs,longtable,fancyhdr,listings,amsmath}
\geometry{margin=2.25cm}
\emergencystretch=2em
\definecolor{azul}{HTML}{17456B}
\hypersetup{colorlinks=true,urlcolor=azul,linkcolor=azul}
\lstdefinestyle{codigo}{
    basicstyle=\ttfamily\footnotesize,
    breaklines=true,
    frame=single,
    columns=fullflexible,
    showstringspaces=false,
    upquote=true,
    keepspaces=true,
    tabsize=4
}
\lstset{style=codigo}
\pagestyle{fancy}\fancyhf{}\lhead{OPP225 | Manual de trabajo}\rhead{Entrega 01}\cfoot{\thepage}
\setlength{\headheight}{14pt}
\setlength{\parskip}{5pt}\setlength{\parindent}{0pt}

\begin{document}

\begin{center}{\LARGE\bfseries Manual semanal de investigación aplicada}\\[4pt]{\large De la roca al clima: fuentes geológicas y transporte potencial de hierro en la Región del Maule}\\[8pt]\end{center}

\section{Parte A. Instalar y configurar las herramientas (2--3 h)}

\subsection*{A1. Conceptos básicos}

\begin{itemize}

    \item \textbf{Python:} Lenguaje de programación que
    utilizaremos para escribir instrucciones, procesar
    datos y realizar análisis.

    \item \textbf{Visual Studio Code (VS Code):} Editor
    de código donde escribiremos, organizaremos y
    ejecutaremos nuestros programas en Python.

    \item \textbf{Git:} Herramienta que permite registrar
    y administrar las versiones de un proyecto.

    Debe instalarse por separado.

    \item \textbf{Terminal:} Interfaz que permite
    comunicarse con el computador mediante comandos.

    \item \textbf{Repositorio:} Carpeta de un proyecto
    cuyo historial de cambios es administrado por Git.

    \item \textbf{CSV:} Formato de archivo que almacena
    datos en filas y columnas, y puede abrirse
    mediante Excel o Python.

\end{itemize}

\subsection*{A2. Instalar y comprobar las herramientas}

Instalar desde las páginas oficiales: \url{https://www.python.org/downloads/}, \url{https://code.visualstudio.com/} y \url{https://git-scm.com/downloads}. En Windows, activar la opción de añadir Python al PATH si el instalador la ofrece. En Ubuntu, se puede usar:

\begin{lstlisting}

sudo apt update
sudo apt install python3 python3-venv python3-pip git
\end{lstlisting}

Abrir VS Code y seleccionar \emph{Terminal > New Terminal}. Ejecutar las comprobaciones siguientes (en Windows, si \texttt{python} no existe, probar \texttt{py}):

\begin{lstlisting}

python --version
git --version
\end{lstlisting}

En Ubuntu usar \texttt{python3 --version}. Registrar las versiones en el informe de entrega. En VS Code, abrir \emph{Extensions} (Ctrl+Shift+X), buscar \textbf{Python} de Microsoft e instalarla; opcionalmente instalar \textbf{Pylance}.

\subsection*{A3. Crear la carpeta del proyecto}

Crear una carpeta llamada \texttt{opp225-maule-vientos} en Documentos. En VS Code elegir \emph{File > Open Folder} y abrirla. Crear estas subcarpetas desde el explorador lateral (clic derecho > New Folder):

\begin{lstlisting}

opp225-maule-vientos/
  datos/
  figuras/
  scripts/
  notas/
  README.md
  requirements.txt
  .gitignore
\end{lstlisting}

En \texttt{README.md}, escribir el nombre del proyecto, la pregunta de investigación, la fecha y la fuente de los datos. En \texttt{.gitignore}, escribir:

\begin{lstlisting}

.venv/
__pycache__/
*.pyc
\end{lstlisting}

\subsection*{A4. Crear el entorno virtual e instalar bibliotecas}

En la terminal abierta en la carpeta del proyecto:

\begin{lstlisting}

# Ubuntu / Linux
python3 -m venv .venv
source .venv/bin/activate
python -m pip install requests pandas matplotlib
python -m pip freeze > requirements.txt
\end{lstlisting}

En Windows (PowerShell), en vez de los dos primeros comandos:

\begin{lstlisting}

py -m venv .venv
.venv\Scripts\Activate.ps1
\end{lstlisting}

Si PowerShell impide activar, usar la terminal \emph{Command Prompt} y ejecutar:

\begin{lstlisting}
.venv\Scripts\activate.bat
\end{lstlisting}

En VS Code seleccionar \emph{Ctrl+Shift+P > Python: Select Interpreter} y escoger el Python dentro de \texttt{.venv}. La terminal debe mostrar \texttt{(.venv)}. No instalar paquetes globalmente.

\subsection*{A5. Primer programa}

Crear \texttt{scripts/hola.py} con:

\begin{lstlisting}[language=Python]

print("Hola, investigacion del viento en Maule")
print("2 + 3 =", 2 + 3)
\end{lstlisting}

Ejecutar \texttt{python scripts/hola.py}. Deben aparecer el saludo y el número 5. Si falla, verificar que la terminal esté en la carpeta principal y que el entorno virtual esté activado.

\subsection*{A6. Primeros pasos con Git}

Git conserva versiones \emph{locales}; no es obligatorio crear una cuenta de GitHub esta semana. En la terminal:

\begin{lstlisting}

git config --global user.name "Nombre Apellido"
git config --global user.email "correo@ejemplo.cl"
git init
git status
git add README.md .gitignore requirements.txt scripts/hola.py
git commit -m "Iniciar proyecto y entorno Python"
git log --oneline
\end{lstlisting}

Reemplazar nombre y correo por datos propios. En VS Code, el ícono \emph{Source Control} muestra cambios; \emph{Commit} guarda una versión. No usar \texttt{git add .} sin revisar archivos; no incorporar contraseñas ni carpetas \texttt{.venv}. Al terminar la semana hacer un segundo commit con código, notas y gráficos.

\subsection{Subir un repositorio local a GitHub}

\textbf{GitHub} permite almacenar en Internet los repositorios
creados con Git, facilitando su respaldo y colaboración.

\begin{enumerate}

    \item \textbf{Crear un repositorio en GitHub:}
    ingresar a \url{https://github.com}, iniciar sesión
    y seleccionar \texttt{New repository}. Asignar un nombre,
    elegir \texttt{Private} o \texttt{Public} y presionar
    \texttt{Create repository}. No inicializar con README
    ni otros archivos.

    \item \textbf{Preparar el repositorio local:}
    abrir la carpeta del proyecto en VS Code y ejecutar
    en la terminal (si todavía no se ha inicializado Git):

    \begin{verbatim}

git init
git add .
git commit -m "Primera version"
\end{verbatim}

    Antes de usar \texttt{git add .}, verificar que
    \texttt{.gitignore} excluya archivos innecesarios
    y credenciales.

    \item \textbf{Conectar Git con GitHub:}
    copiar la dirección HTTPS del repositorio creado
    y ejecutar:

    \begin{verbatim}

git branch -M main
git remote add origin https://github.com/USUARIO/REPOSITORIO.git
\end{verbatim}

    Reemplazar \texttt{USUARIO} y \texttt{REPOSITORIO}
    por los nombres correspondientes.

    \item \textbf{Subir los archivos:}

    \begin{verbatim}

git push -u origin main
\end{verbatim}

    Si se solicita autenticación, seguir las
    instrucciones de GitHub.

    \item \textbf{Subir modificaciones posteriores:}
    después de modificar y guardar los archivos,
    ejecutar:

    \begin{verbatim}

git add .
git commit -m "Descripcion de los cambios"
git push
\end{verbatim}

\end{enumerate}

\noindent
\textbf{Importante:} \texttt{git add} prepara los cambios,
\texttt{git commit} registra una versión local y
\texttt{git push} la envía a GitHub.

\textbf{Verificación:} abrir el repositorio en GitHub,
actualizar la página y comprobar que aparecen los archivos
del proyecto.

\section{Parte B. Obtener datos reales de viento (2--3 h)}

\subsection*{B1. Por qué ERA5 y qué se descargará}

ERA5 es un \textbf{reanálisis}: combina observaciones y un modelo meteorológico para reconstruir condiciones atmosféricas. No es una medición directa en cada punto. Para facilitar el aprendizaje se utilizará la API histórica de \textbf{Open-Meteo}, seleccionando explícitamente \texttt{models=era5}; posteriormente se podrá descargar la grilla ERA5 original de Copernicus para cartografía regional. No mezclar el modelo automático \emph{Best Match} con ERA5 en una comparación temporal.

\textbf{Período piloto:} 1 al 30 de septiembre de 2025, a resolución horaria. \textbf{Puntos exploratorios aproximados:} costa de Constitución ($-35.33,-72.42$), valle de Talca ($-35.43,-71.67$) y precordillera de San Clemente ($-35.55,-71.10$). Las coordenadas son puntos de muestreo iniciales, \emph{no} ubicaciones comprobadas de fuentes de Fe ni receptores definitivos. La extracción se hace por punto de grilla cercano, por lo que no representa microclimas locales.

Se descargarán velocidad (m/s) y dirección (grados meteorológicos, \emph{de donde viene} el viento) a 10 m. El intervalo se expresa en UTC para evitar ambigüedades horarias. La velocidad es una magnitud no negativa; 0 grados indica viento desde el norte y 270 grados desde el oeste.

\subsection*{B2. Crear el programa de descarga}

Copiar el siguiente código \emph{sin modificar nombres de carpetas} en \texttt{scripts/descargar\_viento.py}:

\begin{lstlisting}[language=Python]

from pathlib import Path
import requests
import pandas as pd
BASE = Path(__file__).resolve().parents[1]
DESTINO = BASE / "datos"
DESTINO.mkdir(exist_ok=True)
URL = "https://archive-api.open-meteo.com/v1/archive"
PUNTOS = {
    "Constitucion": (-35.33, -72.42),
    "Talca": (-35.43, -71.67),
    "San_Clemente": (-35.55, -71.10),
}
for nombre, (lat, lon) in PUNTOS.items():
    parametros = {
        "latitude": lat,
        "longitude": lon,
        "start_date": "2025-09-01",
        "end_date": "2025-09-30",
        "hourly": "wind_speed_10m,wind_direction_10m",
        "wind_speed_unit": "ms",
        "timezone": "GMT",
        "models": "era5",
    }
    respuesta = requests.get(URL, params=parametros, timeout=90)
    respuesta.raise_for_status()
    contenido = respuesta.json()
    if "hourly" not in contenido:
        raise ValueError(f"Respuesta inesperada: {contenido}")
    tabla = pd.DataFrame(contenido["hourly"])
    tabla["fecha_hora_utc"] = pd.to_datetime(tabla["time"], utc=True)
    tabla["sitio"] = nombre
    tabla["latitud_solicitada"] = lat
    tabla["longitud_solicitada"] = lon
    tabla["latitud_grilla"] = contenido["latitude"]
    tabla["longitud_grilla"] = contenido["longitude"]
    tabla.to_csv(DESTINO / f"viento_{nombre}.csv", index=False)
    print(nombre, len(tabla), "filas;", "grilla:",
          contenido["latitude"], contenido["longitude"])
\end{lstlisting}

Ejecutar:

\begin{lstlisting}

python scripts/descargar_viento.py
\end{lstlisting}

\textbf{Comprobación:} deben aparecer tres archivos CSV en \texttt{datos/}, cada uno con 720 registros horarios (30 días por 24 horas), si el servicio devuelve todas las horas solicitadas. Si falla la conexión, comprobar internet, esperar y volver a ejecutar. Si aparece un error HTTP, registrar el mensaje; no inventar datos. Conservar los archivos originales y anotar la fecha de descarga.

\section{Parte C. Analizar y graficar (1--2 h)}

Crear \texttt{scripts/analizar\_viento.py} y copiar:

\begin{lstlisting}[language=Python]

from pathlib import Path
import pandas as pd
import matplotlib.pyplot as plt
BASE = Path(__file__).resolve().parents[1]
FIG = BASE / "figuras"
FIG.mkdir(exist_ok=True)
archivos = sorted((BASE / "datos").glob("viento_*.csv"))
if len(archivos) != 3:
    raise RuntimeError("Primero ejecute descargar_viento.py")
df = pd.concat([pd.read_csv(p) for p in archivos], ignore_index=True)
df["fecha_hora_utc"] = pd.to_datetime(df["fecha_hora_utc"], utc=True)
vel = "wind_speed_10m"
dir_ = "wind_direction_10m"
if df[[vel, dir_]].isna().any().any():
    raise ValueError("Hay datos faltantes: revisar CSV")
print("Resumen de velocidades en m/s:")
print(df.groupby("sitio")[vel].agg(["count", "mean", "max"]))
resumen = df.groupby("sitio")[vel].agg(["count", "mean", "max"])
resumen.to_csv(BASE / "datos" / "resumen_viento.csv")
for sitio, parte in df.groupby("sitio"):
    plt.plot(parte["fecha_hora_utc"], parte[vel], label=sitio,
             linewidth=0.7, alpha=0.8)
plt.ylabel("Velocidad a 10 m (m/s)")
plt.xlabel("Fecha (UTC)")
plt.title("Viento horario ERA5, septiembre 2025")
plt.legend()
plt.gcf().autofmt_xdate()
plt.tight_layout()
plt.savefig(FIG / "serie_temporal.png", dpi=160)
plt.close()
# Sectores de PROCEDENCIA del viento, no de destino.
sectores = ["N", "NE", "E", "SE", "S", "SO", "O", "NO"]
fig, axes = plt.subplots(1, 3, figsize=(12, 3.5), sharey=True)
for ax, (sitio, parte) in zip(axes, df.groupby("sitio")):
    indice = (((parte[dir_] % 360) + 22.5) // 45).astype(int) % 8
    frecuencia = indice.value_counts(normalize=True).reindex(range(8), fill_value=0)
    ax.bar(sectores, 100 * frecuencia)
    ax.set_title(sitio)
    ax.set_xlabel("Viento desde")
axes[0].set_ylabel("Frecuencia (%)")
fig.suptitle("Procedencia del viento ERA5, septiembre 2025")
fig.tight_layout()
fig.savefig(FIG / "direcciones.png", dpi=160)
plt.close(fig)
print("Listo: ver carpeta figuras")
\end{lstlisting}

Ejecutar \path{python scripts/analizar_viento.py}. Abrir los dos PNG desde el explorador de VS Code. El primer gráfico muestra la velocidad hora a hora; el segundo muestra desde qué sector proviene el viento, con ocho categorías. \textbf{Importante:} no interpretar los ocho sectores como ocho direcciones de transporte; el aire se desplaza aproximadamente hacia el sector opuesto.

\section{Parte D. Interpretación científica y tareas de reflexión (1 h)}

En \texttt{notas/semana01.md}, responder con palabras propias:

\begin{enumerate}[leftmargin=*]

\item ¿Qué sitio tiene mayor velocidad media y cuál tiene el máximo horario más alto? Indicar valores y unidades obtenidos del CSV, no estimaciones visuales.

\item ¿De qué sector \emph{proviene} con mayor frecuencia el viento en cada sitio? ¿Hacia qué sector se desplazaría una partícula idealmente transportada por ese viento?

\item ¿Cuáles son tres diferencias entre reanálisis, imágenes satelitales y estaciones meteorológicas? ¿Por qué no se puede afirmar que el viento observado transportó hierro?

\item Proponer una hipótesis \emph{comprobable} sobre conectividad atmosférica entre una fuente geológica potencial y un receptor (costa, embalse u otro), aclarando que aún no se ha verificado la localización de la fuente.

\item Identificar al menos tres fuentes de incertidumbre: resolución espacial, altura del viento, topografía, disponibilidad de sedimentos, humedad del suelo o dirección variable.

\end{enumerate}

\textbf{Criterio de calidad:} no inferir emisiones, flujos de Fe ni deposición a partir de la velocidad sola. El umbral de levantamiento de polvo depende de tamaño de grano, humedad, cobertura, rugosidad y esfuerzo cortante; no fijar un umbral universal sin respaldo bibliográfico.

\section{Qué análisis seguiremos desarrollando}

\begin{longtable}{p{0.24\textwidth}p{0.38\textwidth}p{0.24\textwidth}}

\toprule\textbf{Análisis} & \textbf{Pregunta científica} & \textbf{Momento sugerido}\\\midrule
Velocidad y sectores & ¿Cuándo y desde dónde sopla el viento? & Esta semana\\

Frecuencia de eventos intensos & ¿Con qué frecuencia se superan umbrales exploratorios? Comparar varios umbrales, no atribuir erosión automática. & Próxima semana\\

Variación estacional y diurna & ¿Cambia el régimen entre verano/invierno o día/noche? & Noviembre\\

Campo espacial $u,v$ ERA5 & ¿Dónde hay patrones regionales coherentes y cómo afectan los Andes y el litoral? & Noviembre\\

Conectividad fuente--receptor & ¿Los vientos apuntan desde fuentes geológicas cartografiadas hacia receptores? & Tras disponer de capas geológicas\\

Sensibilidad e incertidumbre & ¿Cambian las conclusiones con punto, período, umbral o resolución? & Integración final\\\bottomrule

\end{longtable}

Para el análisis vectorial posterior se emplearán las componentes zonal $u$ (positiva hacia el este) y meridional $v$ (positiva hacia el norte), con $V=\sqrt{u^2+v^2}$. El transporte efectivo requeriría además información sobre disponibilidad y propiedades de partículas, turbulencia, levantamiento y deposición.

\section{Lista de verificación para el viernes 30 de octubre}

\begin{itemize}[label=--,leftmargin=*]

\item VS Code, Python y Git instalados; versiones anotadas.

\item Entorno \texttt{.venv} funcional; \texttt{requirements.txt} presente.

\item Primer programa ejecutado correctamente.

\item Repositorio Git con al menos dos commits; mostrar \texttt{git log --oneline}.

\item Tres CSV horarios y \texttt{resumen\_viento.csv} en \texttt{datos/}.

\item Dos gráficos PNG en \texttt{figuras/}.

\item Respuestas en \texttt{notas/semana01.md}, con una hipótesis y limitaciones.

\item README con fuente, fechas, coordenadas y comandos para reproducir.

\end{itemize}

\section*{Reproducir el proyecto en otro computador}

Instalar Python y Git, copiar la carpeta del proyecto y, desde su raíz, crear/activar un entorno virtual nuevo; ejecutar \texttt{python -m pip install -r requirements.txt}, luego \path{python scripts/descargar_viento.py} y \path{python scripts/analizar_viento.py}. Requiere internet para descargar. Las respuestas pueden cambiar por actualizaciones del servicio: conservar fecha de acceso, CSV originales y versión del entorno. En Git guardar el código, las notas y, si su tamaño lo permite, los CSV de este ejercicio.

\section*{Fuentes y documentación}

\begin{itemize}[leftmargin=*]

\item Copernicus Climate Data Store, \emph{ERA5 hourly data on single levels}: \url{https://cds.climate.copernicus.eu/datasets/reanalysis-era5-single-levels}.

\item Open-Meteo, \emph{Historical Weather API}: \url{https://open-meteo.com/en/docs/historical-weather-api}.

\item Python: \url{https://docs.python.org/3/tutorial/}.

\item VS Code, Python: \url{https://code.visualstudio.com/docs/python/python-tutorial}.

\item Git, libro oficial: \url{https://git-scm.com/book/es/v2}.

\end{itemize}

\textbf{Continuidad editorial:} conservar este archivo como \texttt{manual\_01.tex}. Cada viernes se añadirá una entrega \texttt{manual\_02.tex}, \texttt{manual\_03.tex}, etc., manteniendo la misma estructura: objetivos, procedimiento, código, comprobaciones, interpretación y entregables.

%######################### Lunes 26 de Octubre #################

\section{Descarga de datos meteorológicos ERA5 con Python}

\textbf{Objetivo:} Configurar el acceso a Copernicus y descargar
datos espaciales de viento mediante Python y Jupyter.

\subsection{Crear una cuenta en Copernicus}

\begin{enumerate}

    \item Ingresar a \url{https://cds.climate.copernicus.eu/}
    y crear una cuenta gratuita.

    \item Iniciar sesión y acceder al perfil:
    \url{https://cds.climate.copernicus.eu/profile}.

    \item Buscar la sección \textbf{API key} y copiar
    la clave personal.

    \item Acceder al conjunto de datos
    \href{https://cds.climate.copernicus.eu/datasets/reanalysis-era5-single-levels}{ERA5 hourly data on single levels}
    y aceptar las licencias correspondientes.

\end{enumerate}

\subsection{Configurar la API}

En PowerShell de Windows, ejecutar:

\begin{lstlisting}[style=codigo]

notepad "$HOME\.cdsapirc"
\end{lstlisting}

Escribir las siguientes líneas, reemplazando
\texttt{TU\_CLAVE\_PERSONAL} por la clave obtenida:

\begin{lstlisting}[style=codigo]

url: https://cds.climate.copernicus.eu/api
key: TU_CLAVE_PERSONAL
\end{lstlisting}

Guardar el archivo como \texttt{.cdsapirc},
sin extensión \texttt{.txt}.

\textbf{Importante:} no compartir la clave personal
ni subir este archivo a GitHub.

\subsection{Instalar bibliotecas}

Con el entorno virtual activado, ejecutar en la terminal:

\begin{lstlisting}[style=codigo]

python -m pip install cdsapi xarray netCDF4 numpy pandas matplotlib
\end{lstlisting}

\subsection{Comprobar la configuración desde Jupyter}

Abrir Jupyter en VS Code, seleccionar el entorno virtual
como \textit{kernel} y ejecutar:

\begin{lstlisting}[style=codigo,language=Python]

import cdsapi
cliente = cdsapi.Client()
print("Configuracion de Copernicus correcta.")
\end{lstlisting}

Si no aparecen errores, la configuración se ha leído correctamente.

\subsection{Descargar datos espaciales de viento}

Descargaremos las componentes horizontales del viento
a 10 metros de altura para una zona de la Región del Maule.

\begin{lstlisting}[style=codigo,language=Python]

import cdsapi
cliente = cdsapi.Client()
solicitud = {
    "product_type": ["reanalysis"],
    "variable": [
        "10m_u_component_of_wind",
        "10m_v_component_of_wind"
    ],
    "year": ["2025"],
    "month": ["09"],
    "day": ["01"],
    "time": ["12:00"],
    "area": [-35.0, -72.5, -36.0, -71.0],
    "data_format": "netcdf",
    "download_format": "unarchived"
}
cliente.retrieve(
    "reanalysis-era5-single-levels",
    solicitud,
    "viento_maule.nc"
)
\end{lstlisting}

El parámetro \texttt{area} corresponde a los límites
geográficos en el orden:

\[
[\text{Norte},\text{Oeste},\text{Sur},\text{Este}].
\]

El archivo descargado tendrá formato NetCDF (\texttt{.nc}),
que permite almacenar información espacial y temporal.

\subsection{Visualizar los datos}

En una nueva celda de Jupyter:

\begin{lstlisting}[style=codigo,language=Python]

import xarray as xr
import numpy as np
import matplotlib.pyplot as plt
datos = xr.open_dataset("viento_maule.nc")
print(datos)
u = datos["u10"]
v = datos["v10"]
rapidez = np.sqrt(u**2 + v**2)
rapidez.squeeze().plot()
plt.title("Rapidez del viento - Region del Maule")
plt.show()
\end{lstlisting}

La rapidez del viento se obtiene mediante:

\[
V(x,y,t)=\sqrt{u(x,y,t)^2+v(x,y,t)^2}.
\]

\noindent
\textbf{Resultado esperado:}
Un mapa espacial de rapidez del viento a 10 metros
de altura, expresada en metros por segundo.

\medskip

\noindent
\textbf{Nota:} Si aparece el error
\texttt{403 required licences not accepted},
se deben aceptar las licencias del conjunto ERA5
en el portal de Copernicus.

\clearpage
\section{Entrega 02. Eventos de viento intenso y textura del suelo}

\subsection*{Objetivos}
\begin{enumerate}[label=\arabic*.]
\item Calcular la frecuencia de horas de viento relativamente intenso en tres puntos del Maule.
\item Reconocer el número de eventos consecutivos que superan varios umbrales exploratorios.
\item Obtener y representar espacialmente las fracciones de arena, limo y arcilla del suelo superficial.
\item Distinguir entre características del suelo, condiciones meteorológicas y transporte efectivo de polvo o hierro.
\end{enumerate}

\textbf{Advertencia científica:} los valores de 6, 8 y 10 m/s son umbrales de comparación \emph{exploratorios}, no velocidades críticas de erosión. SoilGrids estima \emph{textura de suelo}, no mediciones de granulometría de aerosoles ni contenido de hierro.

\subsection*{Antes de comenzar}
\texttt{datos/viento\_Constitucion.csv}, \texttt{datos/viento\_Talca.csv} y
\texttt{datos/viento\_San\_Clemente.csv}, creados en la entrega anterior.
Cada uno debiera contener 720 horas de septiembre de 2025. Verificar también
que el entorno \texttt{.venv} esté seleccionado en VS Code.

Instalar la única biblioteca adicional para mapas raster y actualizar el registro:
\begin{lstlisting}
python -m pip install rasterio
python -m pip freeze > requirements.txt
\end{lstlisting}

\subsection{Parte E. Identificar eventos de viento intenso}

\textbf{Pregunta:} ¿en cuál de los tres puntos se presentan más horas con
vientos sobre determinados umbrales y cómo se distribuyen durante el mes?

La frecuencia horaria de excedencia de un umbral \(U\) se calcula como
\[
F(U)=100\,\frac{\#\{t:\,V(t)\geq U\}}{N}\quad(\%).
\]
Un \emph{evento} es una secuencia de una o más horas consecutivas con
\(V(t)\geq U\); dos secuencias separadas por al menos una hora bajo el umbral
se consideran eventos diferentes. No confundir \emph{horas de excedencia}
con \emph{número de eventos}.

Crear el archivo \texttt{scripts/eventos\_viento.py} con el siguiente código:
\begin{lstlisting}[language=Python]
from pathlib import Path
import pandas as pd
import matplotlib.pyplot as plt

BASE = Path(__file__).resolve().parents[1]
DATOS = BASE / 'datos'
FIG = BASE / 'figuras'
FIG.mkdir(exist_ok=True)
archivos = sorted(DATOS.glob('viento_*.csv'))
if len(archivos) != 3:
    raise FileNotFoundError('Se requieren los tres CSV del Manual 01.')

df = pd.concat([pd.read_csv(p) for p in archivos], ignore_index=True)
df['fecha_hora_utc'] = pd.to_datetime(df['fecha_hora_utc'], utc=True)
vel = 'wind_speed_10m'
if df[['sitio', 'fecha_hora_utc', vel]].isna().any().any():
    raise ValueError('Hay valores ausentes en los datos.')
if df.duplicated(['sitio', 'fecha_hora_utc']).any():
    raise ValueError('Hay horas duplicadas por sitio.')

umbrales = [6, 8, 10]  # m/s; exploratorios, no umbrales de erosion
resumen = []
for sitio, grupo in df.groupby('sitio'):
    grupo = grupo.sort_values('fecha_hora_utc').copy()
    dt = grupo['fecha_hora_utc'].diff().dt.total_seconds().div(3600)
    if not (dt.iloc[1:] == 1).all():
        raise ValueError(f'La serie de {sitio} no es horaria continua.')
    for umbral in umbrales:
        excede = grupo[vel] >= umbral
        # Los eventos comienzan cuando una hora supera el umbral
        # y la hora anterior no lo supera.
        inicia = excede & ~excede.shift(fill_value=False)
        resumen.append({
            'sitio': sitio,
            'umbral_m_s': umbral,
            'horas_excedencia': int(excede.sum()),
            'porcentaje_horas': round(100 * excede.mean(), 2),
            'numero_eventos': int(inicia.sum())
        })
res = pd.DataFrame(resumen)
res.to_csv(DATOS / 'resumen_eventos.csv', index=False)
print(res.to_string(index=False))

fig, ax = plt.subplots(figsize=(8, 4))
for sitio, grupo in res.groupby('sitio'):
    ax.plot(grupo['umbral_m_s'], grupo['porcentaje_horas'], 'o-', label=sitio)
ax.set(xlabel='Umbral exploratorio (m/s)',
       ylabel='Horas sobre el umbral (%)',
       title='Frecuencia de vientos intensos, septiembre 2025')
ax.set_xticks(umbrales)
ax.legend()
ax.grid(alpha=0.3)
fig.tight_layout()
fig.savefig(FIG / 'frecuencia_umbrales.png', dpi=160)
plt.close(fig)

# Porcentaje de horas por dia, para el umbral de 8 m/s
copia = df.copy()
copia['dia_utc'] = copia['fecha_hora_utc'].dt.date
copia['excede_8'] = copia[vel] >= 8
diario = (copia.groupby(['sitio', 'dia_utc'])['excede_8']
          .mean().mul(100).reset_index(name='porcentaje_horas'))
diario.to_csv(DATOS / 'eventos_diarios_8ms.csv', index=False)
fig, ax = plt.subplots(figsize=(10, 4))
for sitio, grupo in diario.groupby('sitio'):
    ax.plot(pd.to_datetime(grupo['dia_utc']), grupo['porcentaje_horas'],
            label=sitio, marker='.', linewidth=1)
ax.set(xlabel='Fecha UTC', ylabel='Horas del dia >= 8 m/s (%)',
       title='Variacion diaria de eventos intensos')
ax.legend()
fig.autofmt_xdate()
fig.tight_layout()
fig.savefig(FIG / 'eventos_diarios.png', dpi=160)
plt.close(fig)
print('Listo: CSV de eventos y dos figuras.')

\end{lstlisting}

Ejecutar desde la carpeta principal del proyecto:
\begin{lstlisting}
python scripts/eventos_viento.py
\end{lstlisting}

\textbf{Archivos esperados:}
\begin{itemize}[nosep]
\item \texttt{datos/resumen\_eventos.csv}: horas, porcentajes y número de eventos por sitio y umbral.
\item \texttt{datos/eventos\_diarios\_8ms.csv}: porcentajes diarios para 8 m/s.
\item \texttt{figuras/frecuencia\_umbrales.png} y \texttt{figuras/eventos\_diarios.png}.
\end{itemize}

\textbf{Comprobación:} cada porcentaje debe estar entre 0 y 100; para cada sitio,
la frecuencia de excedencia debe disminuir o permanecer constante cuando el
umbral aumenta. Interpretar el resultado como comparación meteorológica,
no como una tasa de emisión de partículas.

\subsection{Parte F. Mapa de textura superficial del suelo (25 min)}

\textbf{Pregunta:} ¿cómo varía espacialmente la proporción estimada de arena,
limo y arcilla en un sector piloto cercano a Talca?

\textbf{Fuente real:} \emph{SoilGrids}, ISRIC -- World Soil Information,
\url{https://isric.org/explore/soilgrids}. Se utilizarán las capas
\texttt{sand\_0-5cm\_mean}, \texttt{silt\_0-5cm\_mean} y
\texttt{clay\_0-5cm\_mean}, correspondientes a estimaciones para el
intervalo 0--5 cm de profundidad. Resolución nominal: aproximadamente
250 m. Licencia: CC BY 4.0.

Estas variables representan fracciones de la tierra fina (no incluyen toda la
fracción de fragmentos gruesos). Los datos se publican en g/kg:
\[
\mathrm{porcentaje}\;(\%)=\frac{\mathrm{valor\ en\ g/kg}}{10}.
\]
Una distribución de estas tres fracciones \emph{no} equivale a una curva
granulométrica completa; no permite calcular D10, D50 o D90.

Para evitar descargar mapas de todo Chile, se usará un \textbf{rectángulo
piloto alrededor de Talca}, entre longitudes $-71.85$ y $-71.45$ y
latitudes $-35.65$ y $-35.20$. No debe rotularse como ``toda la Región del Maule''.

Crear la subcarpeta \texttt{datos/soilgrids/} y el archivo
\texttt{scripts/textura\_suelo.py}. Este código consulta el servicio WCS,
guarda tres GeoTIFF, produce mapas y genera una tabla descriptiva:
\begin{lstlisting}[language=Python]
from pathlib import Path
import requests
import numpy as np
import pandas as pd
import matplotlib.pyplot as plt
import rasterio
from rasterio.plot import plotting_extent

BASE = Path(__file__).resolve().parents[1]
DATOS = BASE / 'datos' / 'soilgrids'
FIG = BASE / 'figuras'
DATOS.mkdir(parents=True, exist_ok=True)
FIG.mkdir(exist_ok=True)

# Sector piloto de Talca (NO equivale a toda la Region del Maule).
OESTE, ESTE = -71.85, -71.45
SUR, NORTE = -35.65, -35.20
URL = 'https://maps.isric.org/mapserv'

# Obtener rasteres de textura 0-5 cm mediante WCS.
for propiedad in ('sand', 'silt', 'clay'):
    destino = DATOS / f'{propiedad}_0-5cm_mean.tif'
    if destino.exists():
        print('Usando archivo existente:', destino)
        continue
    parametros = [
        ('map', f'/map/{propiedad}.map'),
        ('SERVICE', 'WCS'),
        ('VERSION', '2.0.1'),
        ('REQUEST', 'GetCoverage'),
        ('COVERAGEID', f'{propiedad}_0-5cm_mean'),
        ('FORMAT', 'GEOTIFF_INT16'),
        ('SUBSET', f'long({OESTE},{ESTE})'),
        ('SUBSET', f'lat({SUR},{NORTE})'),
        ('SUBSETTINGCRS', 'http://www.opengis.net/def/crs/EPSG/0/4326'),
        ('OUTPUTCRS', 'http://www.opengis.net/def/crs/EPSG/0/4326')
    ]
    try:
        r = requests.get(URL, params=parametros, timeout=90)
        r.raise_for_status()
        if r.content[:4] not in (b'II*\x00', b'MM\x00*'):
            raise ValueError('El servidor no devolvio un GeoTIFF. ' +
                             r.text[:180].replace('\n', ' '))
        destino.write_bytes(r.content)
        print('Descargado:', destino)
    except (requests.RequestException, ValueError) as error:
        raise RuntimeError(
            'Fallo WCS: ' + str(error) + '\n' +
            'Alternativa: descargar las capas sand/silt/clay de 0-5 cm '
            'en SoilGrids.org o mediante QGIS-WCS, guardarlas con los '
            'nombres indicados en datos/soilgrids y repetir este script.') from error

capas = {}
referencia = None
for propiedad in ('sand', 'silt', 'clay'):
    with rasterio.open(DATOS / f'{propiedad}_0-5cm_mean.tif') as src:
        if referencia is None:
            referencia = (src.crs, src.transform, src.width, src.height)
            extent = plotting_extent(src)
        elif (src.crs, src.transform, src.width, src.height) != referencia:
            raise ValueError('Rasteres con grillas distintas: alinear en QGIS.')
        a = src.read(1, masked=True).astype('float64')
        # SoilGrids usa g/kg; dividiendo por 10 se obtienen porcentajes.
        capas[propiedad] = a / 10.0

fig, axes = plt.subplots(1, 3, figsize=(12, 4), constrained_layout=True)
for ax, (prop, nombre) in zip(axes, [('sand','Arena'),
                                      ('silt','Limo'), ('clay','Arcilla')]):
    im = ax.imshow(capas[prop], extent=extent, origin='upper',
                   cmap='viridis', vmin=0, vmax=100)
    ax.set(title=f'{nombre}, suelo 0-5 cm (%)',
           xlabel='Longitud', ylabel='Latitud')
    fig.colorbar(im, ax=ax, shrink=0.7)
fig.suptitle('SoilGrids: textura estimada, sector de Talca')
fig.savefig(FIG / 'textura_talca.png', dpi=160)
plt.close(fig)

# Estadisticos espaciales descriptivos; celdas modeladas, NO muestras.
filas = []
for prop, a in capas.items():
    valores = a.compressed()
    if not len(valores):
        raise ValueError(f'Sin datos validos para {prop}.')
    filas.append({'componente': prop, 'celdas_validas': len(valores),
                  'media_pct': round(float(np.mean(valores)), 2),
                  'min_pct': round(float(np.min(valores)), 2),
                  'max_pct': round(float(np.max(valores)), 2)})
res = pd.DataFrame(filas)
res.to_csv(DATOS / 'resumen_textura.csv', index=False)
print(res.to_string(index=False))
print('Listo: mapa y tabla de textura estimada.')

\end{lstlisting}

Ejecutar desde el directorio principal del proyecto:
\begin{lstlisting}
python scripts/textura_suelo.py
\end{lstlisting}

\textbf{Archivos esperados:} tres archivos GeoTIFF en
\texttt{datos/soilgrids/}; la figura \texttt{figuras/textura\_talca.png} y
la tabla \texttt{datos/soilgrids/resumen\_textura.csv}.

\textbf{Comprobaciones:} los porcentajes deben estar en el intervalo
0--100 y las tres capas deben cubrir exactamente la misma grilla.
Los porcentajes de las tres fracciones pueden no sumar exactamente 100
por tratarse de predicciones modeladas de forma independiente.
No interpretar diferencias entre píxeles como mediciones de terreno.

\subsubsection*{Si SoilGrids no responde}
ISRIC informa interrupciones de su \emph{API REST}; el procedimiento anterior
usa \emph{WCS}, que es un servicio distinto, pero también requiere Internet y
puede fallar. \textbf{No generar números ficticios ni reemplazar las capas
por imágenes sin datos.} Como alternativa:
\begin{enumerate}[nosep]
\item Abrir \url{https://soilgrids.org/} y comprobar si su interfaz permite
exportar el recorte de \texttt{sand}, \texttt{silt} y \texttt{clay}, 0--5 cm, promedio.
\item Si se dispone de QGIS, crear una conexión WCS con la URL
\url{https://maps.isric.org/mapserv?map=/map/sand.map}; repetir cambiando
\texttt{sand} por \texttt{silt} y \texttt{clay}. Exportar la capa correspondiente
a GeoTIFF y utilizar \emph{exactamente} la misma extensión, resolución y
sistema de referencia para las tres exportaciones.
\item Guardar los archivos con los nombres indicados en el código, dentro de
\texttt{datos/soilgrids/}, y volver a ejecutar el script.
\item Si ninguno de los servicios está disponible, registrar la incidencia
en \texttt{notas/semana02.md} y completar la interpretación meteorológica.
La parte espacial queda pendiente de obtener datos reales, sin inventarlos.
\end{enumerate}

\textbf{Nota:} las descargas por WCS dependen del servidor y de sus parámetros.
La URL y la documentación oficial están en
\url{https://docs.isric.org/globaldata/soilgrids/wcs.html}.

\subsection{Parte G. Interpretación y entrega (10 min de discusión)}
Crear \texttt{notas/semana02.md} y responder:
\begin{enumerate}[label=\arabic*.]
\item ¿Qué sitio presenta mayor porcentaje de horas sobre 8 m/s?
¿Coincide con el sitio que presenta más eventos? Justificar con la tabla.
\item ¿Cómo cambia la frecuencia al pasar de 6 a 10 m/s?
¿Cuáles son las limitaciones de usar un solo mes y tres puntos?
\item En el sector analizado, ¿dónde predomina la arena, el limo o la arcilla?
Apoyar la respuesta en los mapas y en la tabla, si se obtuvieron los GeoTIFF.
\item ¿Por qué \emph{no} es correcto concluir que los sectores con más limo
emiten necesariamente más polvo, o que el viento intenso transportó hierro?
\item Formular una hipótesis comprobable que conecte textura del suelo,
intensidad del viento y una posible fuente geológica, indicando qué datos
adicionales harían falta para evaluarla (humedad, cobertura, Fe, etc.).
\end{enumerate}

\textbf{Entrega mínima:} código ejecutable; dos gráficos meteorológicos;
\texttt{resumen\_eventos.csv}; respuestas documentadas; repositorio actualizado.
\textbf{Entrega adicional, condicionada a disponibilidad del servidor:}
tres raster de textura, mapa y tabla de porcentajes. Registrar fecha,
fuente, recorte geográfico y posibles errores de descarga.

\textbf{Antes del viernes 13 de noviembre:} completar las figuras, revisar
la variación diurna de los eventos y buscar una fuente regional de
\emph{granulometría medida en sedimentos o suelos} y evidencia mineralógica
de hierro. SoilGrids es un punto de partida espacial, no reemplaza datos
geológicos de la fuente de polvo.

\subsection*{Fuentes de la Entrega 02}
\begin{itemize}[nosep]
\item ISRIC, \emph{SoilGrids}: \url{https://isric.org/explore/soilgrids}.
\item ISRIC, WCS: \url{https://docs.isric.org/globaldata/soilgrids/wcs.html}.
\item ISRIC, unidades de las capas: \url{https://docs.isric.org/globaldata/soilgrids/SoilGrids_faqs_01.html}.
\item Copernicus, ERA5: \url{https://cds.climate.copernicus.eu/datasets/reanalysis-era5-single-levels}.
\end{itemize}



\section{Parte H. Cartografía geológica y contenido de hierro}
\textbf{Objetivo:} incorporar información geológica \emph{real} de la Región
 del Maule, delimitar posibles áreas de interés y distinguir evidencias
 litológicas de mediciones de contenido de hierro (Fe).

\textbf{Importante:} una carta geológica representa unidades y tipos de roca;
\emph{no} es un mapa de concentración de Fe ni permite asignar porcentajes
numéricos de hierro a cada unidad. Una roca que contiene minerales con Fe
no necesariamente genera polvo disponible para transporte. Esta actividad
sirve para formular y priorizar hipótesis, no para cuantificar emisiones.

\subsection{H1. Fuentes oficiales que utilizaremos}
\begin{enumerate}[label=\arabic*.,leftmargin=*]
\item \textbf{Mapa geológico de Chile 1:1.000.000.} El Ministerio del Medio
Ambiente (MMA), con información de SERNAGEOMIN, publica \emph{Geología Base}
en formatos SHP y GeoJSON. Incluye ubicación y atributos geológicos;
 su escala es de reconocimiento regional, no de detalle local:
\url{https://lineasdebasepublicas.mma.gob.cl/datos_abiertos/dataset/geologia-base}.
\item \textbf{Cartas de mayor detalle.} En el catálogo de SERNAGEOMIN buscar
\emph{Geología del área San Clemente--Melado} (2024, código M433)
y \emph{Geología del área Pichibelco--Cauquenes} (2023, código M434),
ambas a escala 1:100.000:
\url{https://portalsdnmext.sernageomin.cl/publicacionesgeologicas}.
Estas cartas deben emplearse para interpretar las unidades presentes en
sectores concretos; su cobertura no es toda la región.
\item \textbf{Geoquímica de sedimentos.} SERNAGEOMIN cuenta con un programa
de cartografía geoquímica y un visor de sedimentos:
\url{https://www.sernageomin.cl/geologia-general/}.
Si existen muestras georreferenciadas en el área de estudio con análisis de
Fe, sus valores pueden servir para contrastar las hipótesis cartográficas.
\end{enumerate}

\subsection{H2. Descargar la cartografía en formato GeoJSON (15 min)}
\begin{enumerate}[nosep]
\item Acceder a la página \emph{Geología Base} indicada en H1.
\item Bajo \emph{Datos y recursos}, elegir \emph{Geología Base GeoJSON}
y luego \emph{Descargar}. El recurso se ofrece como archivo ZIP.
\item Crear la carpeta \texttt{datos/geologia/} en el proyecto.
\item Descomprimir el ZIP dentro de esa carpeta, de manera que el archivo
\texttt{.geojson} quede allí o en una subcarpeta.
\item Registrar la URL, fecha de acceso, nombre exacto del archivo y
\emph{escala 1:1.000.000} en \texttt{notas/semana02.md}.
\end{enumerate}

\textbf{Si falla la descarga:} abrir el catálogo oficial y consultar
las cartas M433/M434 para la interpretación. No crear geometrías ni
unidades geológicas ficticias. El código siguiente deberá ejecutarse
cuando el archivo real esté disponible.

\subsection{H3. Crear un mapa geológico regional con Python}
En la terminal del entorno virtual instalar:
\begin{lstlisting}
python -m pip install geopandas
\end{lstlisting}
Crear \texttt{scripts/mapa\_geologico\_maule.py} con el siguiente código.
 El programa busca el GeoJSON descargado, selecciona una \textbf{ventana
geográfica exploratoria} del Maule, muestra las unidades y crea un CSV
con las unidades cartografiadas. \textbf{La ventana es un rectángulo de
análisis, no el límite administrativo regional.}

\begin{lstlisting}[language=Python]
"""Mapa geologico exploratorio del Maule (datos reales desde MMA/SERNAGEOMIN).
Requiere GeoJSON descargado desde el portal MMA, sin inventar Fe.
"""
from pathlib import Path
import geopandas as gpd
import matplotlib.pyplot as plt
import pandas as pd

BASE = Path(__file__).resolve().parents[1]
ORIGEN = BASE / 'datos' / 'geologia'
FIG = BASE / 'figuras'
FIG.mkdir(parents=True, exist_ok=True)
ORIGEN.mkdir(parents=True, exist_ok=True)
archivos = list(ORIGEN.rglob('*.geojson')) + list(ORIGEN.rglob('*.json'))
if not archivos:
    raise FileNotFoundError(
        'Descargue y descomprima Geologia Base GeoJSON desde el portal MMA '
        'en datos/geologia/. No se usan datos de ejemplo como resultados.')

# El GeoJSON de geologia base es nacional; mostramos una ventana del Maule.
gdf = gpd.read_file(archivos[0]).to_crs('EPSG:4326')
gdf = gdf[gdf.geometry.notna() & ~gdf.geometry.is_empty].copy()
# Recorte por interseccion con el rectangulo (no equivale al limite regional).
oeste, este, sur, norte = -73.4, -70.3, -36.5, -34.4
maule = gdf.cx[oeste:este, sur:norte].copy()
if maule.empty:
    raise ValueError('No se encontraron unidades en el recorte del Maule. Revisar archivo/CRS.')

# Detectar columna de descripcion de rocas; no asumir un esquema fijo.
columnas = [c for c in maule.columns if c != 'geometry']
preferencias = ['tipo_roca', 'tipo de roca', 'litologia', 'litologia',
                'composicio', 'descripcion', 'descripcio', 'geo']
col = next((c for p in preferencias for c in columnas if p in c.lower()), None)
if col is None:
    col = next((c for c in columnas if pd.api.types.is_string_dtype(maule[c])), None)
if col is None:
    raise ValueError(f'No hay campo textual para clasificar unidades. Columnas: {columnas}')

maule['unidad_mapa'] = maule[col].fillna('Sin clasificar').astype(str).str.strip()
conteos = maule['unidad_mapa'].value_counts()
principales = set(conteos.head(8).index)
maule['clase'] = maule['unidad_mapa'].where(maule['unidad_mapa'].isin(principales),
                                          'Otras unidades')
fig, ax = plt.subplots(figsize=(9, 7))
maule.plot(column='clase', categorical=True, legend=True, ax=ax,
           edgecolor='0.45', linewidth=0.25,
           legend_kwds={'loc': 'upper left', 'bbox_to_anchor': (1.02, 1),
                        'fontsize': 7})
ax.set_xlim(oeste, este)
ax.set_ylim(sur, norte)
ax.set_xlabel('Longitud (grados)')
ax.set_ylabel('Latitud (grados)')
ax.set_title('Geologia base: recorte exploratorio del Maule\n'
             'Fuente: MMA / SERNAGEOMIN, escala 1:1.000.000')
fig.tight_layout()
fig.savefig(FIG / 'geologia_maule.png', dpi=180, bbox_inches='tight')
plt.close(fig)

# Listado de unidades representadas, NO porcentajes de Fe.
resumen = (maule.groupby('unidad_mapa', dropna=False).size()
           .rename('poligonos_en_recorte').reset_index()
           .sort_values('poligonos_en_recorte', ascending=False))
resumen.to_csv(ORIGEN / 'unidades_maule.csv', index=False, encoding='utf-8-sig')
print(f'GeoJSON: {archivos[0].name}')
print(f'Campo cartografico empleado: {col}')
print(f'Poligonos en el recorte: {len(maule)}')
print('Archivos: figuras/geologia_maule.png, datos/geologia/unidades_maule.csv')
print('ATENCION: este mapa NO contiene concentraciones ni porcentajes de Fe.')
\end{lstlisting}
Ejecutar desde la carpeta principal:
\begin{lstlisting}
python scripts/mapa_geologico_maule.py
\end{lstlisting}
\textbf{Resultado:} \texttt{figuras/geologia\_maule.png} y
\texttt{datos/geologia/unidades\_maule.csv}. La columna usada para
agrupar unidades se imprime en pantalla. Comprobar que sus categorías
sean realmente tipos o descripciones geológicas; si el archivo emplea
otro esquema, revisar sus atributos antes de interpretar el mapa.

\subsection{H4. ¿Cómo reconocer posibles sectores de interés para Fe?}
Revisar la leyenda y la memoria de las cartas detalladas. Buscar, si
están presentes, unidades máficas, rocas volcánicas con minerales
ferromagnesianos, cuerpos mineralizados o depósitos con minerales
portadores de hierro. \textbf{Estas son pistas cualitativas, no medidas
de Fe ni garantías de sedimentos erosionables.} Los depósitos
superficiales y aluviales podrían ser importantes para la disponibilidad
de partículas, pero requieren verificación de textura, humedad, cobertura
y mineralogía.

Para documentar candidatos, crear en \texttt{notas/semana02.md} una
tabla con \emph{sector, unidad/código, litología descrita, argumento
mineralógico, evidencia de sedimento superficial, fuente documental y
grado de incertidumbre}. Registrar al menos dos sectores si la
cartografía los permite. No clasificarlos automáticamente como fuentes
emisoras de hierro.

\subsection{H5. ¿Es posible calcular un porcentaje de hierro?}
\textbf{Sí, pero únicamente si existe un análisis geoquímico asociado
a una muestra y una fracción/material definidos.} Si el dato está
expresado como concentración másica de hierro en partes por millón:
\[
\mathrm{Fe\ (\%)}=\frac{\mathrm{Fe\ (ppm)}}{10\,000}.
\]
Si se informa \emph{hierro total como} $\mathrm{Fe_2O_3}$ (porcentaje
en masa), la masa equivalente de Fe se calcula, si la convención analítica
lo permite, como
\[
\mathrm{Fe\ (\%)}=\mathrm{Fe_2O_3\ (\%)}
\frac{2M_{\mathrm{Fe}}}{2M_{\mathrm{Fe}}+3M_{\mathrm{O}}}
\simeq 0.6994\,\mathrm{Fe_2O_3\ (\%)}.
\]
Este factor \textbf{no} significa que toda la muestra contenga ese
óxido: es una conversión de la forma de reportar el elemento.
No aplicar conversiones sin revisar las unidades, el método y si el
valor corresponde a Fe total o a alguna fracción extraíble.

\textbf{Ejemplo exclusivamente aritmético, no un resultado del Maule:}
$25\,000$ ppm de Fe equivalen a $2.5$ por ciento en masa. No utilizar
este valor en mapas o conclusiones del proyecto.

\subsection{H6. Mapa de Fe medido (actividad condicionada a datos reales)}
Cuando se obtenga una tabla geoquímica verificable con coordenadas,
Fe y la fuente de cada registro, guardar
\texttt{datos/geologia/muestras\_fe.csv} con las columnas
\texttt{latitud,longitud,fe\_pct,fuente}. Para preparar el CSV,
transformar las unidades de la fuente según H5 y dejar anotado
\emph{método, material analizado, tamaño de fracción y calidad del dato}.
\textbf{No se proporciona un CSV ficticio.}

Crear \texttt{scripts/mapa\_fe\_maule.py}:
\begin{lstlisting}[language=Python]
"""Mapa opcional de Fe medido: solo ejecutable con una tabla real y trazable."""
from pathlib import Path
import pandas as pd
import matplotlib.pyplot as plt

BASE = Path(__file__).resolve().parents[1]
archivo = BASE / 'datos' / 'geologia' / 'muestras_fe.csv'
if not archivo.exists():
    raise FileNotFoundError(
        'No hay datos geoquimicos medidos. No crear porcentajes ficticios. '
        'Cuando consiga una tabla verificable, guardela como datos/geologia/muestras_fe.csv')

df = pd.read_csv(archivo)
obligatorias = {'latitud', 'longitud', 'fe_pct', 'fuente'}
if not obligatorias.issubset(df.columns):
    raise ValueError(f'Faltan columnas: {sorted(obligatorias - set(df.columns))}')
for c in ['latitud', 'longitud', 'fe_pct']:
    df[c] = pd.to_numeric(df[c], errors='raise')
if (df.fe_pct < 0).any() or (df.fe_pct > 100).any():
    raise ValueError('Fe (%) debe estar entre 0 y 100; comprobar unidades.')
if df[['latitud','longitud','fe_pct']].isna().any().any():
    raise ValueError('Hay valores faltantes.')
fig, ax = plt.subplots(figsize=(8, 6))
p = ax.scatter(df.longitud, df.latitud, c=df.fe_pct, s=65,
               cmap='viridis', edgecolors='black', linewidths=0.4)
fig.colorbar(p, ax=ax, label='Fe medido (% en masa)')
ax.set(xlabel='Longitud', ylabel='Latitud',
       title='Concentracion de hierro en muestras geoquimicas')
fig.tight_layout()
destino = BASE / 'figuras' / 'muestras_fe_maule.png'
destino.parent.mkdir(exist_ok=True)
fig.savefig(destino, dpi=170)
print('Figura creada:', destino)
print('Verificar: muestras, fraccion granulometrica, metodo y fuente.')
\end{lstlisting}
El programa solo genera \texttt{figuras/muestras\_fe\_maule.png} si
se dispone de las muestras. Representa puntos medidos y \emph{no}
interpola las concentraciones entre ellos.

\subsection{H7. Actividad de interpretación (para completar después de la clase)}
\begin{enumerate}[label=\arabic*.]
\item Seleccionar dos unidades geológicas del mapa y describir su
composición litológica, citando una carta o memoria geológica.
\item ¿Qué información faltaría para afirmar que una unidad contiene
más Fe que otra? Diferenciar Fe en roca, suelo y sedimento fino.
\item ¿Qué lugares del mapa parecen interesantes para contrastar con
la frecuencia de eventos de viento intenso? Justificar, sin afirmar
transporte real.
\item ¿Por qué el mapa regional a escala 1:1.000.000 no identifica
por sí solo afloramientos expuestos, granulometría o áreas de emisión?
\end{enumerate}

\textbf{Entregable adicional:} mapa geológico y tabla de unidades,
\emph{si se pudo descargar el GeoJSON}; dos sectores candidatos
justificados con cartografía oficial. El mapa de Fe es una actividad
posterior \emph{condicionada} a disponer de análisis químicos.

\textbf{Nota para la próxima sesión:} revisar las cartas M433 y M434,
la disponibilidad real de datos de Fe en el visor geoquímico y la
coincidencia entre fuentes superficiales posibles, textura y vientos.

\end{document}

